\documentclass[11pt,a4paper]{article}
% ============================================================
% preambulo.tex — preámbulo común de la documentación teórica
% de SIL_Flight_Control (docs/teoria/).
%
% Uso, desde un documento en docs/teoria/<tema>/:
%   \documentclass[11pt,a4paper]{article}
%   % ============================================================
% preambulo.tex — preámbulo común de la documentación teórica
% de SIL_Flight_Control (docs/teoria/).
%
% Uso, desde un documento en docs/teoria/<tema>/:
%   \documentclass[11pt,a4paper]{article}
%   % ============================================================
% preambulo.tex — preámbulo común de la documentación teórica
% de SIL_Flight_Control (docs/teoria/).
%
% Uso, desde un documento en docs/teoria/<tema>/:
%   \documentclass[11pt,a4paper]{article}
%   \input{../comun/preambulo}
% ============================================================

% ---------- Idioma y codificación ----------
\usepackage[utf8]{inputenc}
\usepackage[T1]{fontenc}
% es-noshorthands: no redefine < > " (rompen TikZ y siunitx)
% es-nodecimaldot: el punto decimal se queda como punto, igual que en el código
\usepackage[spanish,es-noshorthands,es-nodecimaldot,es-tabla]{babel}
\usepackage{lmodern}
\usepackage{microtype}

% ---------- Página ----------
\usepackage[margin=2.5cm]{geometry}
\usepackage{parskip}

% ---------- Matemáticas y unidades ----------
\usepackage{amsmath,amssymb}
\usepackage{bm}
\usepackage[output-decimal-marker={.},per-mode=symbol]{siunitx}

% ---------- Tablas, figuras, código ----------
\usepackage{booktabs}
\usepackage{array}
\usepackage{graphicx}
\usepackage{tikz}
\usetikzlibrary{arrows.meta,positioning}
\usepackage{listings}
\lstset{
  basicstyle=\ttfamily\small,
  columns=fullflexible,
  keepspaces=true,
  frame=single,
  framerule=0.3pt,
  breaklines=true,
  inputencoding=utf8,
  extendedchars=true,
  literate={á}{{\'a}}1 {é}{{\'e}}1 {í}{{\'i}}1 {ó}{{\'o}}1 {ú}{{\'u}}1
           {ñ}{{\~n}}1 {Δ}{{$\Delta$}}1 {°}{{$^\circ$}}1,
}

% ---------- Referencias cruzadas ----------
\usepackage[colorlinks=true,linkcolor=blue!50!black,citecolor=green!40!black,
            urlcolor=blue!60!black]{hyperref}

% ---------- Macros del proyecto ----------
% Nombres de código: archivos, funciones, claves de configuración y columnas
\newcommand{\archivo}[1]{\texttt{#1}}
\newcommand{\func}[1]{\texttt{#1}}
\newcommand{\clave}[1]{\texttt{#1}}
\newcommand{\proyecto}{\textsf{SIL\_Flight\_Control}}

% Notación
\newcommand{\dpm}{\Delta p}            % presión diferencial medida
\newcommand{\rhoo}{\rho_0}             % densidad a nivel del mar ISA
\newcommand{\Normal}{\mathcal{N}}      % distribución normal
\newcommand{\E}{\mathbb{E}}            % esperanza
\DeclareMathOperator{\atantwo}{atan2}

% ============================================================

% ---------- Idioma y codificación ----------
\usepackage[utf8]{inputenc}
\usepackage[T1]{fontenc}
% es-noshorthands: no redefine < > " (rompen TikZ y siunitx)
% es-nodecimaldot: el punto decimal se queda como punto, igual que en el código
\usepackage[spanish,es-noshorthands,es-nodecimaldot,es-tabla]{babel}
\usepackage{lmodern}
\usepackage{microtype}

% ---------- Página ----------
\usepackage[margin=2.5cm]{geometry}
\usepackage{parskip}

% ---------- Matemáticas y unidades ----------
\usepackage{amsmath,amssymb}
\usepackage{bm}
\usepackage[output-decimal-marker={.},per-mode=symbol]{siunitx}

% ---------- Tablas, figuras, código ----------
\usepackage{booktabs}
\usepackage{array}
\usepackage{graphicx}
\usepackage{tikz}
\usetikzlibrary{arrows.meta,positioning}
\usepackage{listings}
\lstset{
  basicstyle=\ttfamily\small,
  columns=fullflexible,
  keepspaces=true,
  frame=single,
  framerule=0.3pt,
  breaklines=true,
  inputencoding=utf8,
  extendedchars=true,
  literate={á}{{\'a}}1 {é}{{\'e}}1 {í}{{\'i}}1 {ó}{{\'o}}1 {ú}{{\'u}}1
           {ñ}{{\~n}}1 {Δ}{{$\Delta$}}1 {°}{{$^\circ$}}1,
}

% ---------- Referencias cruzadas ----------
\usepackage[colorlinks=true,linkcolor=blue!50!black,citecolor=green!40!black,
            urlcolor=blue!60!black]{hyperref}

% ---------- Macros del proyecto ----------
% Nombres de código: archivos, funciones, claves de configuración y columnas
\newcommand{\archivo}[1]{\texttt{#1}}
\newcommand{\func}[1]{\texttt{#1}}
\newcommand{\clave}[1]{\texttt{#1}}
\newcommand{\proyecto}{\textsf{SIL\_Flight\_Control}}

% Notación
\newcommand{\dpm}{\Delta p}            % presión diferencial medida
\newcommand{\rhoo}{\rho_0}             % densidad a nivel del mar ISA
\newcommand{\Normal}{\mathcal{N}}      % distribución normal
\newcommand{\E}{\mathbb{E}}            % esperanza
\DeclareMathOperator{\atantwo}{atan2}

% ============================================================

% ---------- Idioma y codificación ----------
\usepackage[utf8]{inputenc}
\usepackage[T1]{fontenc}
% es-noshorthands: no redefine < > " (rompen TikZ y siunitx)
% es-nodecimaldot: el punto decimal se queda como punto, igual que en el código
\usepackage[spanish,es-noshorthands,es-nodecimaldot,es-tabla]{babel}
\usepackage{lmodern}
\usepackage{microtype}

% ---------- Página ----------
\usepackage[margin=2.5cm]{geometry}
\usepackage{parskip}

% ---------- Matemáticas y unidades ----------
\usepackage{amsmath,amssymb}
\usepackage{bm}
\usepackage[output-decimal-marker={.},per-mode=symbol]{siunitx}

% ---------- Tablas, figuras, código ----------
\usepackage{booktabs}
\usepackage{array}
\usepackage{graphicx}
\usepackage{tikz}
\usetikzlibrary{arrows.meta,positioning}
\usepackage{listings}
\lstset{
  basicstyle=\ttfamily\small,
  columns=fullflexible,
  keepspaces=true,
  frame=single,
  framerule=0.3pt,
  breaklines=true,
  inputencoding=utf8,
  extendedchars=true,
  literate={á}{{\'a}}1 {é}{{\'e}}1 {í}{{\'i}}1 {ó}{{\'o}}1 {ú}{{\'u}}1
           {ñ}{{\~n}}1 {Δ}{{$\Delta$}}1 {°}{{$^\circ$}}1,
}

% ---------- Referencias cruzadas ----------
\usepackage[colorlinks=true,linkcolor=blue!50!black,citecolor=green!40!black,
            urlcolor=blue!60!black]{hyperref}

% ---------- Macros del proyecto ----------
% Nombres de código: archivos, funciones, claves de configuración y columnas
\newcommand{\archivo}[1]{\texttt{#1}}
\newcommand{\func}[1]{\texttt{#1}}
\newcommand{\clave}[1]{\texttt{#1}}
\newcommand{\proyecto}{\textsf{SIL\_Flight\_Control}}

% Notación
\newcommand{\dpm}{\Delta p}            % presión diferencial medida
\newcommand{\rhoo}{\rho_0}             % densidad a nivel del mar ISA
\newcommand{\Normal}{\mathcal{N}}      % distribución normal
\newcommand{\E}{\mathbb{E}}            % esperanza
\DeclareMathOperator{\atantwo}{atan2}


\title{Control lateral-direccional\\[0.3em]
       \large Yaw damper, alabeo y rumbo en cascada para el Ugly Stick}
\author{Proyecto \proyecto}
\date{Octubre de 2026 --- Paso 4 del plan, fase F2 (diseño)}

% Gráficas: datos de dynamics/control_lat_design.py --datos datos/
\pgfplotsset{
  diseno/.style={width=0.92\linewidth, height=5.2cm, grid=major,
                 grid style={gray!25}, tick label style={font=\small},
                 label style={font=\small}, legend style={font=\footnotesize}},
}

\begin{document}
\maketitle

\begin{abstract}
Se diseñan las leyes de control lateral-direccional del simulador: un
\emph{yaw damper} con \emph{washout} sobre el timón, un lazo de alabeo
($p$, $\phi$) sobre los alerones y un lazo externo de rumbo que comanda el
ángulo de alabeo. El diseño es clásico, por cierre sucesivo de lazos, sobre el
mismo modelo lineal que integra el simulador en C, con los servos incluidos.
Con las ganancias elegidas el \emph{dutch roll} pasa de $\zeta = 0.10$ a
$0.83$, la espiral queda estable, el alabeo responde en \SI{0.53}{s} con 2\%
de sobrepaso y el rumbo se establece en \SI{7.6}{s} sin sobrepaso. Todos los
polos complejos tienen $\zeta \ge 0.42$; los márgenes de ganancia son
$\ge \SI{22.8}{dB}$ y el lazo tolera \SI{160}{ms} de retardo. El documento
explica cada decisión, en particular un acoplamiento entre el lazo de rumbo y
el \emph{dutch roll} que obligó a subir la ganancia del yaw damper.
\end{abstract}

\tableofcontents

% ============================================================
\section{Objetivo y alcance}
% ============================================================

El Ugly Stick en lazo abierto tiene dos problemas laterales (tabla~\ref{tab:la}):
el \emph{dutch roll} está poco amortiguado ($\zeta = 0.10$, por debajo del mínimo
de nivel~1 de MIL-F-8785C~\cite{mil8785}, $\zeta \ge 0.19$) y la espiral es
inestable ($T_2 = \SI{9.3}{s}$). El objetivo del Paso~4 es un piloto automático
lateral con tres modos, seleccionados con el canal \clave{mode\_switch} del radio:

\begin{center}
\begin{tabular}{cll}
\toprule
\clave{mode\_switch} & Modo & Lazos activos \\
\midrule
0 & Manual        & ninguno (stick directo a las superficies) \\
1 & Estabilizado  & yaw damper + alabeo; el stick de roll comanda $\phi_c$ \\
2 & Rumbo         & yaw damper + alabeo + rumbo; $\psi_c$ se captura al entrar \\
\bottomrule
\end{tabular}
\end{center}

Este documento cubre el \textbf{diseño} (fase F2). La implementación en C es
de las fases F3--F6; la sección~\ref{sec:impl} fija lo que el C debe respetar.
Fuera de alcance: control longitudinal y estimación de actitud (el control
toma $\phi$ y $\psi$ de un AHRS ideal: estado verdadero más ruido).

% ============================================================
\section{Criterios de diseño}
% ============================================================

Acordados antes de diseñar; \archivo{dynamics/control\_lat\_design.py} los
verifica y sale con error si alguno no se cumple:

\begin{center}
\small
\begin{tabular}{p{0.46\linewidth}lp{0.27\linewidth}}
\toprule
Criterio & Valor & Motivo \\
\midrule
$\zeta$ de todo polo complejo en lazo cerrado & $\ge 0.4$ & MIL nivel 1 ($0.19$) con margen \\
Espiral & estable & modos 1 y 2 \\
Escalón de $\phi$: sobrepaso / subida 10--90\% & $\le 5\%$ / $\le \SI{0.8}{s}$ & lazo interno rápido y limpio \\
Escalón de $\psi$: sobrepaso / establecimiento 5\% & $\le 5\%$ / $\le \SI{10}{s}$ & sin pasarse del rumbo \\
Margen de ganancia (subida) por canal & $\ge \SI{6}{dB}$ & incertidumbre de derivadas \\
Retardo puro tolerado en ambos canales & $\ge \SI{50}{ms}$ & lazo a \SI{200}{Hz} con holgura \\
\bottomrule
\end{tabular}
\end{center}

% ============================================================
\section{Modelo de diseño}
% ============================================================

\subsection{Planta}

Se usa el modelo lateral del simulador (\func{build\_lateral\_matrices},
\archivo{src/dynamics.c}), en ejes cuerpo, con el rumbo como quinto estado:
\begin{equation}
  \dot{\bm x} = A\bm x + B\bm u, \qquad
  \bm x = [\,v,\ p,\ r,\ \phi,\ \psi\,]^T, \quad \bm u = [\,\delta_a,\ \delta_r\,]^T,
\end{equation}
\begin{equation}
A = \begin{bmatrix}
Y_v & Y_p + w_0 & -(u_0 - Y_r) & g\cos\theta_0 & 0\\
L_v & L_p & L_r & 0 & 0\\
N_v & N_p & N_r & 0 & 0\\
0 & 1 & \tan\theta_0 & 0 & 0\\
0 & 0 & \sec\theta_0 & 0 & 0
\end{bmatrix},\qquad
B = \begin{bmatrix}
0 & Y_{\delta r}\\ L_{\delta a} & L_{\delta r}\\ N_{\delta a} & N_{\delta r}\\ 0 & 0\\ 0 & 0
\end{bmatrix}.
\end{equation}
Las derivadas se leen de \archivo{aircraft/ugly\_stick.json}, el mismo archivo
que carga el simulador; un test verifica que el bloque $4\times4$ coincide con
la matriz del C. A la planta se le agregan los \textbf{servos} como sistemas de
primer orden ($\tau_a = \SI{0.05}{s}$, $\tau_r = \SI{0.07}{s}$, del mismo JSON)
y el estado del \emph{washout} (sección~\ref{sec:yd}): 8 estados en total.

\subsection{Lazo abierto}

La tabla~\ref{tab:la} resume los modos de la planta sin control. El roll es
rápido y bien amortiguado; los problemas son el dutch roll y la espiral. El
quinto eigenvalor, en 0, es el rumbo: $\psi$ no realimenta a nada, solo
integra $r$.

\begin{table}[!htb]
\centering
\caption{Modos laterales del Ugly Stick en lazo abierto (sin servos).}
\label{tab:la}
\begin{tabular}{lccc}
\toprule
Modo & $\lambda$ & $\omega_n$ [\si{rad/s}] & $\zeta$ \\
\midrule
Roll           & $-13.31$            & ---  & --- \\
Dutch roll     & $-0.413 \pm 4.096j$ & 4.12 & 0.10 \\
Espiral        & $+0.075$            & ---  & inestable, $T_2 = \SI{9.3}{s}$ \\
Rumbo          & $0$                 & ---  & integrador ($\psi$) \\
\bottomrule
\end{tabular}
\end{table}

\subsection{Convención de signos}

Convención de Nelson~\cite{nelson}, la misma del simulador:
$\delta_a > 0$ produce alabeo a la derecha ($L_{\delta a} = 109.9 > 0$) y
$\delta_r > 0$ (borde de salida a la izquierda) produce guiñada a la izquierda
($N_{\delta r} = -9.69 < 0$). Por eso las tres leyes de abajo tienen ganancias
\textbf{positivas}:
\begin{align}
  \delta_r &= K_r\, r_{wo}, \label{eq:yd}\\
  \delta_a &= K_\phi(\phi_c - \phi) - K_p\, p, \label{eq:roll}\\
  \phi_c   &= \operatorname{sat}_{\pm\phi_{\max}}\!\bigl(K_\psi(\psi_c - \psi)\bigr). \label{eq:hdg}
\end{align}
Un $r > 0$ (nariz a la derecha) da $\delta_r > 0$, que produce un momento de
guiñada negativo y se opone al movimiento: eso es amortiguar.

% ============================================================
\section{Yaw damper}
\label{sec:yd}
% ============================================================

\subsection{Ley y washout}

El timón realimenta la tasa de guiñada \emph{lavada}:
\begin{equation}
  r_{wo} = \frac{\tau_w s}{1 + \tau_w s}\, r,
  \qquad \tau_w = \SI{1}{s}.
\end{equation}
El \emph{washout} es un filtro pasa-altas: deja pasar la oscilación del dutch
roll ($\omega \approx \SI{4}{rad/s}$, muy por encima de $1/\tau_w$) y bloquea
la componente constante de $r$. Sin él, en un viraje sostenido
($r \ne 0$ estacionario) el yaw damper metería timón contra el viraje
\cite{stevens}. En el modelo de diseño es un estado más,
$\dot x_{wo} = (r - x_{wo})/\tau_w$ con $r_{wo} = r - x_{wo}$.

\subsection{Elección de \texorpdfstring{$K_r$}{Kr}}

La figura~\ref{fig:lugar} muestra los polos complejos al variar $K_r$ de 0 a
1.2. El dutch roll se amortigua muy rápido: $K_r = 0.20$ ya da $\zeta = 0.37$ y
$K_r \approx 0.27$ da $\zeta = 0.5$. Con el yaw damper solo, la espiral sigue
inestable ($+0.069$): el timón no la corrige, eso le toca al lazo de alabeo.

\begin{figure}[htbp]
\centering
\begin{tikzpicture}
\begin{axis}[diseno, width=0.8\linewidth, xlabel={Parte real [\si{1/s}]}, ylabel={Parte imaginaria [\si{rad/s}]},
             xmin=-14, xmax=1, ymin=0, ymax=12,
             colorbar, colormap/viridis, point meta min=0, point meta max=1.2,
             colorbar style={ylabel={$K_r$}, ylabel style={font=\small}}]
  \addplot[scatter, only marks, mark=*, mark size=1pt, scatter src=explicit]
    table[x=re, y=im, meta=Kr] {datos/lugar_yaw_damper.dat};
  \addplot[only marks, mark=x, mark size=4pt, very thick, red]
    coordinates {(-0.413,4.096)};
  \addplot[only marks, mark=o, mark size=4pt, very thick, red]
    coordinates {(-5.488,3.758) (-2.317,0.554)};
  \draw[dashed, gray] (axis cs:0,0) -- (axis cs:-13.75,12) node[pos=0.55, above, font=\scriptsize, sloped] {$\zeta=0.75$};
  \draw[dashed, gray] (axis cs:0,0) -- (axis cs:-5.24,12) node[pos=0.55, above, font=\scriptsize, sloped] {$\zeta=0.4$};
\end{axis}
\end{tikzpicture}
\caption{Lugar de las raíces del yaw damper (polos complejos, servos y washout
incluidos). $\times$: dutch roll en lazo abierto; $\circ$: $K_r = 0.42$
(dutch roll en $-5.49 \pm 3.76j$ y un polo lento del washout acoplado en
$-2.32 \pm 0.55j$).}
\label{fig:lugar}
\end{figure}

El valor final, $K_r = 0.42$, es mayor que el que bastaría para el yaw damper
solo. La razón está en el lazo de rumbo (sección~\ref{sec:acople}). Con
$K_r = 0.42$ el dutch roll queda en $-5.49 \pm 3.76j$ ($\zeta = 0.83$,
$\omega_n = \SI{6.65}{rad/s}$).

% ============================================================
\section{Alabeo}
\label{sec:alabeo}
% ============================================================

\subsection{Estructura}

La ley~\eqref{eq:roll} son dos lazos anidados en una sola expresión: $K_p$
amortigua el alabeo (realimentación de $p$ del giróscopo) y $K_\phi$ cierra la
actitud. Como $\phi$ integra a $p$, el lazo de $\phi$ es de tipo~1 respecto a
$\delta_a$ y no necesita integrador para seguir un escalón.

\subsection{Elección de \texorpdfstring{$K_p$ y $K_\phi$}{Kp y Kphi}}

Con el yaw damper cerrado se barrieron $K_p \in \{0.05, 0.10, 0.15\}$ y
$K_\phi \in [0.3, 1.6]$ midiendo la respuesta al escalón de $\phi_c$:

\begin{center}
\begin{tabular}{cccccc}
\toprule
$K_p$ [\si{s}] & $K_\phi$ & $t_r$ [\si{s}] & sobrepaso & $t_{s,5\%}$ [\si{s}] \\
\midrule
0.05 & 0.8 & 0.33 & 3.2\% & 0.54 \\
0.10 & 0.5 & 0.90 & 3.0\% & 1.19 \\
\textbf{0.10} & \textbf{0.8} & \textbf{0.55} & \textbf{2.8\%} & \textbf{0.79} \\
0.10 & 1.2 & 0.27 & 2.5\% & 0.48 \\
0.15 & 0.8 & 0.69 & 2.4\% & 0.94 \\
\bottomrule
\end{tabular}
\end{center}

Se eligió $K_p = 0.10$, $K_\phi = 0.8$: cumple con holgura y no es tan agresivo
como para saturar los alerones con escalones normales. Con $K_\phi = 0.8$, un
escalón de \SI{30}{\degree} pide de entrada
$0.8 \times \SI{0.52}{rad} \approx \SI{24}{\degree}$ de alerón y satura brevemente en
\SI{20}{\degree}; uno de \SI{10}{\degree} pide \SI{5.2}{\degree} como máximo
(figura~\ref{fig:phi}). Un $K_p$ más alto vuelve oscilatorio el par
servo--modo de roll ($K_p = 0.2$: $-16.5 \pm 20.6j$, $\zeta = 0.62$).

Con el modo estabilizado (yaw damper + alabeo) todos los polos son estables,
\textbf{incluida la espiral}, y $\zeta_{\min} = 0.76$. El estado
estacionario de $\phi$ queda 1.8\% por encima de $\phi_c$: en un viraje
sostenido $L_r\,r \ne 0$ necesita algo de alerón, que sale de un pequeño error.

\begin{figure}[htbp]
\centering
\begin{tikzpicture}
\begin{axis}[diseno, xlabel={$t$ [\si{s}]}, ylabel={[\si{\degree}], [\si{\degree/s}]},
             xmin=0, xmax=4, legend pos=north east]
  \addplot[thick, blue] table[x=t, y=phi_deg] {datos/escalon_phi.dat};
  \addlegendentry{$\phi$}
  \addplot[thick, red, dashed] table[x=t, y=da_deg] {datos/escalon_phi.dat};
  \addlegendentry{$\delta_a$}
  \addplot[thick, green!50!black, dotted] table[x=t, y=r_dps] {datos/escalon_phi.dat};
  \addlegendentry{$r$}
  \addplot[gray, densely dashed] coordinates {(0,10) (4,10)};
\end{axis}
\end{tikzpicture}
\caption{Escalón de $\phi_c = \SI{10}{\degree}$ en modo estabilizado (modelo
lineal con servos).}
\label{fig:phi}
\end{figure}

% ============================================================
\section{Rumbo}
\label{sec:rumbo}
% ============================================================

\subsection{Ley}

En un viraje coordinado $\dot\psi = (g/V)\tan\phi \approx (g/V)\,\phi$
\cite{beard}, así que un lazo proporcional~\eqref{eq:hdg} que comanda $\phi_c$
da, si el lazo de alabeo es mucho más rápido, una respuesta de primer orden
con constante $\tau_\psi \approx V/(g K_\psi) = \SI{2.6}{s}$ para $K_\psi = 1$.
La regla de separación de anchos de banda~\cite{beard} pide que el lazo externo
sea al menos 5 veces más lento que el interno; el de alabeo tiene
$t_r = \SI{0.53}{s}$, así que $K_\psi \approx 1$ está en el límite razonable.
El error de rumbo se envuelve a $[-\pi, \pi)$ y $\phi_c$ se limita a
$\pm\phi_{\max} = \pm\SI{30}{\degree}$.

\subsection{Acoplamiento con el dutch roll}
\label{sec:acople}

Al cerrar el lazo de rumbo sobre el diseño inicial ($K_r = 0.27$) apareció un
par de polos cerca de \SI{2.9}{rad/s} cuyo amortiguamiento \emph{cae} al
subir $K_\psi$, hasta volverse inestable en $K_\psi \approx 2$
(figura~\ref{fig:acople}). Es el dutch roll ya amortiguado: el rumbo se
corrige con alerón, el alerón tiene guiñada adversa
($N_{\delta a} = -2.51$) y esa guiñada excita el dutch roll a través del
lazo. Con $K_r = 0.27$, para tener $\zeta \ge 0.4$ habría que bajar a
$K_\psi \approx 0.5$ (rumbo en $\sim\SI{15}{s}$).

La solución fue subir el yaw damper: con $K_r = 0.42$ se llega a
$K_\psi = 1.0$ con $\zeta_{\min} = 0.42$. Más $K_r$ (0.60) da más margen en
$\zeta$, pero el dutch roll se parte en dos polos reales y la actividad del
timón crece sin beneficio en la respuesta de rumbo.

\begin{figure}[htbp]
\centering
\begin{tikzpicture}
\begin{axis}[diseno, xlabel={$K_\psi$}, ylabel={$\zeta_{\min}$ en modo rumbo},
             xmin=0.2, xmax=2, ymin=0, ymax=0.8, legend pos=north east]
  \addplot[thick, red] table[x=Kpsi, y=zmin_Kr027] {datos/zeta_rumbo.dat};
  \addlegendentry{$K_r = 0.27$}
  \addplot[thick, blue] table[x=Kpsi, y=zmin_Kr042] {datos/zeta_rumbo.dat};
  \addlegendentry{$K_r = 0.42$ (elegido)}
  \addplot[thick, green!50!black, dashed] table[x=Kpsi, y=zmin_Kr060] {datos/zeta_rumbo.dat};
  \addlegendentry{$K_r = 0.60$}
  \addplot[gray, densely dashed] coordinates {(0.2,0.4) (2,0.4)};
  \addplot[only marks, mark=o, mark size=3pt, blue, thick] coordinates {(1.0,0.42)};
\end{axis}
\end{tikzpicture}
\caption{Amortiguamiento mínimo de los polos complejos en modo rumbo contra
$K_\psi$ ($K_p = 0.10$, $K_\phi = 0.8$). Línea punteada: criterio
$\zeta \ge 0.4$.}
\label{fig:acople}
\end{figure}

\subsection{Respuesta}

Con $K_\psi = 1.0$ (figura~\ref{fig:psi}) un cambio de rumbo de
\SI{30}{\degree} se completa sin sobrepaso: subida en \SI{4.9}{s},
establecimiento al 5\% en \SI{7.6}{s}, sin error estacionario. El alabeo
pedido llega a $\approx\SI{30.5}{\degree}$, justo en el límite de
$\phi_{\max}$. Dos detalles de la figura: al principio $\psi$ baja unos
\SI{3}{\degree}, que es la guiñada adversa del alerón ($N_{\delta a} < 0$)
mientras el avión entra en el viraje; y la meseta entre 2 y \SI{3}{s} es el
modo menos amortiguado del lazo ($-1.06 \pm 2.27j$, $\zeta = 0.42$,
sección~\ref{sec:acople}). Cambios de rumbo mayores quedan limitados por la saturación
(viraje a \SI{30}{\degree} de alabeo, $\dot\psi \approx \SI{13}{\degree/s}$). El
timón se mueve poco ($\le \SI{8.6}{\degree}$, ya que el washout deja pasar solo
el transitorio).

\begin{figure}[htbp]
\centering
\begin{tikzpicture}
\begin{axis}[diseno, xlabel={$t$ [\si{s}]}, ylabel={[\si{\degree}]},
             xmin=0, xmax=15, legend pos=north east]
  \addplot[thick, blue] table[x=t, y=psi_deg] {datos/escalon_psi.dat};
  \addlegendentry{$\psi$}
  \addplot[thick, orange] table[x=t, y=phi_deg] {datos/escalon_psi.dat};
  \addlegendentry{$\phi$}
  \addplot[thick, red, dashed] table[x=t, y=da_deg] {datos/escalon_psi.dat};
  \addlegendentry{$\delta_a$}
  \addplot[thick, green!50!black, dotted] table[x=t, y=dr_deg] {datos/escalon_psi.dat};
  \addlegendentry{$\delta_r$}
  \addplot[gray, densely dashed] coordinates {(0,30) (15,30)};
\end{axis}
\end{tikzpicture}
\caption{Escalón de $\psi_c = \SI{30}{\degree}$ en modo rumbo (modelo lineal,
sin la saturación de $\phi_c$).}
\label{fig:psi}
\end{figure}

% ============================================================
\section{Robustez}
% ============================================================

Se evaluó en el modo rumbo, el más exigente, de dos formas:

\paragraph{Rango de ganancia por canal.} Multiplicando por $k$ la ley de un
canal (con el otro en su valor nominal) y buscando el rango de $k$ estable:

\begin{center}
\begin{tabular}{lccc}
\toprule
Canal & $k$ estable & Margen al subir & Margen al bajar \\
\midrule
Alerón ($K_p$, $K_\phi$, $K_\psi$) & $(0,\ >50]$ & $> \SI{34}{dB}$ & cualquier $k > 0$ \\
Timón ($K_r$)                         & $[0.14,\ 13.8]$ & \SI{22.8}{dB} & \SI{-17.1}{dB} \\
\bottomrule
\end{tabular}
\end{center}

El límite inferior del timón tiene una lectura importante: \textbf{en modo
rumbo el yaw damper es imprescindible}. Con menos del 14\% de $K_r$ el
acoplamiento de la sección~\ref{sec:acople} desestabiliza el dutch roll. Si el
yaw damper falla en vuelo, el modo 2 debe desactivarse.

\paragraph{Retardo puro.} Con un retardo $T$ en ambas entradas (aproximación
de Padé de orden 2) el lazo sigue estable hasta $T = \SI{160}{ms}$. El lazo
en C correrá a \SI{200}{Hz}: el retenedor de orden cero equivale a
$\approx\SI{2.5}{ms}$ de retardo, dos órdenes de magnitud por debajo.

% ============================================================
\section{Implementación prevista en C}
\label{sec:impl}
% ============================================================

Lo que las fases F3--F6 deben respetar para que el C reproduzca este diseño:

\begin{itemize}
  \item \textbf{Frecuencia}: leyes en \func{Task\_5ms} (\SI{200}{Hz}).
  \item \textbf{Entradas}: $p$, $r$ del giróscopo (con ruido); $\phi$, $\psi$
        del AHRS ideal. Unidades en grados o radianes indistintamente: las
        ganancias son adimensionales ($K_p$ en segundos).
  \item \textbf{Washout discreto} (Euler hacia atrás, estable para cualquier
        $\Delta t$):
        $x_{wo,k} = x_{wo,k-1} + \frac{\Delta t}{\tau_w + \Delta t}(r_k - x_{wo,k-1})$,
        $r_{wo,k} = r_k - x_{wo,k}$.
  \item \textbf{Salidas}: $\delta_a$, $\delta_r$ son deflexiones
        \emph{respecto al trim}; el C las pasa a comando de servo en
        \si{\micro s} con $u = \delta/k_{\text{servo}} + \text{trim}_{us}$
        (y signo de Nelson, sin la inversión del pedal, que es solo del piloto).
  \item \textbf{Saturaciones}: $\phi_c \in \pm\SI{30}{\degree}$; las de
        las superficies ($\pm\SI{20}{\degree}$, $\pm\SI{30}{\degree}$) y los
        limitadores de velocidad ya están en los servos.
  \item \textbf{Modo 1}: $\phi_c = \phi_{\max}\cdot$\,stick de roll; el pedal se
        suma al yaw damper como mando directo.
  \item \textbf{Modo 2}: $\psi_c$ se captura al entrar al modo; con el stick
        de roll fuera de una banda muerta el piloto manda $\phi_c$ y al
        soltarlo se recaptura el rumbo. Error de rumbo envuelto a
        $[-\pi,\pi)$.
  \item \textbf{Cambios de modo}: reiniciar el estado del washout al entrar a
        un modo con yaw damper, para evitar escalones de timón.
\end{itemize}

% ============================================================
\section{Resumen y reproducción}
% ============================================================

\begin{center}
\begin{tabular}{llll}
\toprule
Ganancia & Valor & Unidad & Lazo \\
\midrule
$K_r$          & 0.42 & \si{rad/(rad/s)} & yaw damper (timón) \\
$\tau_w$       & 1.0  & \si{s}           & washout \\
$K_p$          & 0.10 & \si{rad/(rad/s)} & amortiguamiento de alabeo (alerón) \\
$K_\phi$       & 0.80 & \si{rad/rad}     & actitud de alabeo (alerón) \\
$K_\psi$       & 1.00 & \si{rad/rad}     & rumbo $\rightarrow \phi_c$ \\
$\phi_{\max}$  & 30   & \si{\degree}     & límite de $\phi_c$ \\
\bottomrule
\end{tabular}
\end{center}

Para reproducir los números y las figuras, desde la raíz del repo:
\begin{lstlisting}
python3 dynamics/control_lat_design.py aircraft/ugly_stick.json \
    --datos docs/teoria/control/datos
\end{lstlisting}
El script escribe \archivo{dynamics/ugly\_control\_lat.json} (ganancias,
polos, respuestas y márgenes) y los \archivo{.dat} de las figuras. Si se
cambian las derivadas del avión o los servos, hay que volver a correrlo: sale
con error si el diseño deja de cumplir los criterios.

% ============================================================
\section{Limitaciones}
% ============================================================

\begin{itemize}
  \item Diseño sobre el \textbf{modelo lineal} en una sola condición de vuelo
        (\SI{25.25}{m/s}, \SI{1000}{m}). Las derivadas escalan con $\bar q$:
        a otra velocidad las ganancias cambian de efecto (sin
        \emph{gain scheduling} por ahora).
  \item Derivadas \textbf{estimadas}, no medidas (TODO 8.12). Los márgenes de
        ganancia son amplios, pero el acoplamiento de la
        sección~\ref{sec:acople} depende de $N_{\delta a}$ y conviene
        revisarlo cuando haya datos de vuelo.
  \item Actitud de un \textbf{AHRS ideal}: un estimador real agrega retardo y
        error en $\phi$ y $\psi$.
  \item Sin turbulencia ni viento: el rumbo es el de la nariz, no el de la
        trayectoria (sin corrección por deriva).
\end{itemize}

% ============================================================
\begin{thebibliography}{9}
% ============================================================
\bibitem{nelson}
R.~C. Nelson, \emph{Flight Stability and Automatic Control}, 2.ª ed.
McGraw-Hill, 1998, caps.~5 y 9.

\bibitem{stevens}
B.~L. Stevens, F.~L. Lewis y E.~N. Johnson, \emph{Aircraft Control and
Simulation}, 3.ª ed. Wiley, 2016, cap.~4 (yaw damper con washout).

\bibitem{beard}
R.~W. Beard y T.~W. McLain, \emph{Small Unmanned Aircraft: Theory and
Practice}. Princeton University Press, 2012, cap.~6 (cierre sucesivo de lazos).

\bibitem{mil8785}
MIL-F-8785C, \emph{Military Specification: Flying Qualities of Piloted
Airplanes}, 1980.
\end{thebibliography}

\end{document}
